% Generated draft: requires mathematical and source review.
\documentclass[11pt]{article}
\usepackage[margin=25mm]{geometry}
\usepackage{fontspec}
\setmainfont{Latin Modern Roman}
\usepackage{amsmath,amssymb,mathtools,unicode-math}
\setmathfont{Latin Modern Math}
\usepackage{enumitem,needspace,xurl,hyperref}
\setlength{\parindent}{0pt}
\setlength{\parskip}{5pt}
\setlength{\emergencystretch}{3em}
\newcommand{\N}{\mathbb{N}}
\newcommand{\Z}{\mathbb{Z}}
\newcommand{\Q}{\mathbb{Q}}
\newcommand{\R}{\mathbb{R}}
\newcommand{\C}{\mathbb{C}}
\newcommand{\F}{\mathbb{F}}
\newcommand{\A}{\mathbb{A}}
\newcommand{\PP}{\mathbb{P}}
\newcommand{\E}{\mathbb{E}}
\newcommand{\eps}{\varepsilon}
\newcommand{\dd}{\,\mathrm{d}}
\begin{document}
\section*{UnsolvedMath ID: 30003740}
\textbf{Source problem code:} catalogue entry
\subsection{Definitions and mathematical statement}
\textbf{Scope and conditional notation.} The source describes a research question about phenotypically structured populations, but does not specify a population model, trait space, or class of environmental fluctuations. The following notation names the requested observables without choosing those missing data. If a model represents the population at time \(t\) by a finite nonnegative measure \(\mu_t\) on a measurable trait space \((X,\mathcal X)\), define its total population size by
\[
N(t)=\mu_t(X).
\]
Whenever \(N(t)>0\), define its normalized phenotype distribution by
\[
p_t(A)=\frac{\mu_t(A)}{N(t)},\qquad A\in\mathcal X.
\]
This measure representation is an optional framework for expressing the question, not a model asserted by the source. Write \(e(t)\) for the environment; in a continuous-time periodic formulation, periodicity means that there exists \(T>0\) such that \(e(t+T)=e(t)\) whenever both times belong to the time domain. The source does not specify the time domain or a mathematical meaning of general environmental fluctuation.

\textbf{Research question.} For phenotypically structured populations in a time-periodic or fluctuating environment, determine the conditions for survival or extinction, describe population-size dynamics, identify which traits are selected, and characterize how environmental variation affects the phenotype distribution. In the conditional notation above, these requests concern \(N(t)\), \(\mu_t\), and, while the population is nonzero, \(p_t\).

\textbf{Unresolved qualifications.} A precise model-dependent statement requires the evolution law, admissible parameters and environments, trait and time domains, initial data, any boundary conditions, and definitions of survival, extinction, and trait selection. The source does not distinguish, for example, eventual disappearance from asymptotic extinction, or survival from persistence bounded away from zero. It also does not specify whether trait selection means concentration, convergence of normalized distributions, or another property. Consequently, no particular evolution equation, limiting assertion, universal quantifier over models, or survival criterion is supplied here.

The supplied dated assessment reports results for particular models in periodic or steadily shifting environments and describes the fully general fluctuating-environment question as open. It does not give hypotheses or theorem statements identifying the scope of those results.
\subsection{Short English statement}
Determine when a phenotypically structured population survives or becomes extinct in a periodic or fluctuating environment, how its size evolves, which traits are selected, and how environmental changes affect its trait distribution. The supplied record leaves the population model and the meanings of these outcomes unspecified, so this remains a research question requiring further hypotheses. Its supplied status is \emph{partially solved}; that status is not independently verified here.
\subsection{Sources}
\begin{itemize}
\item Supplied local registry record: permanent UnsolvedMath ID \(30003740\), title \emph{Trait Selection in Fluctuating Environments}, source problem code \texttt{catalogue entry}, in the Oberwolfach Reports Open Problems collection. The proposed year is recorded as 2018; no individual proposer is supplied.
\item Supplied report citation: \emph{The Mathematics of Mechanobiology and Cell Signaling (2019)}, Oberwolfach Reports, pp.\ 433--506, DOI \texttt{10.4171/owr/2018/8}. Supplied source URL: \url{https://doi.org/10.4171/owr/2018/8}. The background attributes the problem to workshop open-problem discussions documented in Oberwolfach Reports; no exact problem page or speaker is supplied.
\item Supplied literature source labeled ``original OWR report'': \url{https://doi.org/10.4171/OWR/2018/8}.
\item Supplied literature source labeled ``Mirrahimi--Perthame--Souganidis, time fluctuations'': \url{https://doi.org/10.1016/j.anihpc.2013.10.001}.
\item Supplied literature source labeled ``Figueroa Iglesias--Mirrahimi, periodic environments'': \url{https://doi.org/10.1137/18M1175185}.
\item Literature assessment date: \texttt{2026-08-21}. This is supplied dated curation, not an independently verified assessment. It labels the problem partially solved and reports that models for periodic or steadily shifting environments establish survival criteria, population-size behavior, and concentration of phenotypic distributions, while the fully general fluctuating-environment question remains open. The supplied background calls this dated literature triage derived from OWR curation and requests independent verification of time-sensitive status and cited claims.
\end{itemize}
\end{document}
